% =====================================================================
%  SICAK JÜPİTERLERDE YÖRÜNGE DÖNEMİ DEĞİŞİMİ — ADIM ADIM PROJE KILAVUZU
%  Derleme: pdflatex kilavuz.tex  (iki kez çalıştır: içindekiler için)
%  Overleaf: yeni proje -> bu dosyayı yapıştır -> "sekiller" klasörüne PNG'leri yükle
%  Şekil dosyası yoksa belge yine derlenir; yerine gri bir kutu çıkar.
% =====================================================================
\documentclass[11pt,a4paper]{article}
\usepackage[utf8]{inputenc}
\usepackage[T1]{fontenc}
\usepackage{lmodern}
\usepackage[shorthands=off,turkish]{babel}
\usepackage[margin=2.3cm]{geometry}
\usepackage{amsmath,amssymb}
\usepackage{graphicx}
\usepackage{booktabs,longtable,tabularx,array}
\usepackage{enumitem}
\usepackage[table]{xcolor}
\usepackage{caption}
\usepackage{float}
\usepackage{listings}
\usepackage[most]{tcolorbox}
\usepackage{pgfplots}
\pgfplotsset{compat=1.17}
\usepackage{fancyhdr}
\usepackage[hidelinks]{hyperref}
% Ek A'daki düz metinde (Word'e kopyalanacak) geçen özel karakterler
\DeclareUnicodeCharacter{2032}{\ensuremath{'}}
\DeclareUnicodeCharacter{03C3}{\ensuremath{\sigma}}
\DeclareUnicodeCharacter{2212}{\ensuremath{-}}
\DeclareUnicodeCharacter{0394}{\ensuremath{\Delta}}
\DeclareUnicodeCharacter{00B7}{\ensuremath{\cdot}}
\DeclareUnicodeCharacter{00B1}{\ensuremath{\pm}}
\DeclareUnicodeCharacter{00B2}{\ensuremath{^2}}
\DeclareUnicodeCharacter{03C7}{\ensuremath{\chi}}
\DeclareUnicodeCharacter{03C0}{\ensuremath{\pi}}
\DeclareUnicodeCharacter{00BD}{\ensuremath{\tfrac{1}{2}}}
\DeclareUnicodeCharacter{2265}{\ensuremath{\geq}}
\DeclareUnicodeCharacter{2248}{\ensuremath{\approx}}

\setlength{\parskip}{0.45em}
\setlength{\parindent}{0pt}
\renewcommand{\arraystretch}{1.2}
\pagestyle{fancy}\fancyhf{}
\lhead{\small O$-$C Proje Kılavuzu}\rhead{\small TÜBİTAK 2209-A}\cfoot{\thepage}

% ---------- kutular ----------
\newtcolorbox{basit}[1][]{colback=green!5,colframe=green!45!black,title={#1},fonttitle=\bfseries,breakable}
\newtcolorbox{dikkat}[1][]{colback=red!4,colframe=red!60!black,title={#1},fonttitle=\bfseries,breakable}
\newtcolorbox{ornek}[1][]{colback=blue!4,colframe=blue!50!black,title={#1},fonttitle=\bfseries,breakable}
\newtcolorbox{hoca}[1][]{colback=yellow!8,colframe=orange!70!black,title={#1},fonttitle=\bfseries,breakable}

% ---------- şekil: dosya yoksa kutu ----------
\newcommand{\sekil}[2]{\IfFileExists{#1}{\includegraphics[width=#2]{#1}}{\fbox{\parbox[c][4cm][c]{#2}{\centering\small (Şekil dosyası bulunamadı:\\\texttt{\detokenize{#1}}\\Kod bölümündeki betiği çalıştır.)}}}}

% ---------- kod ----------
\definecolor{kodzemin}{rgb}{0.97,0.97,0.97}
\lstset{basicstyle=\ttfamily\footnotesize,backgroundcolor=\color{kodzemin},frame=single,
  breaklines=true,columns=fullflexible,keepspaces=true,showstringspaces=false,
  commentstyle=\color{green!40!black},keywordstyle=\color{blue!70!black},stringstyle=\color{red!60!black},
  language=Python,upquote=true,
  literate={ı}{{\i}}1 {ğ}{{\u{g}}}1 {ö}{{\"o}}1 {ü}{{\"u}}1 {ş}{{\c{s}}}1 {ç}{{\c{c}}}1
           {İ}{{\.I}}1 {Ö}{{\"O}}1 {Ü}{{\"U}}1 {Ş}{{\c{S}}}1 {Ç}{{\c{C}}}1 {Ğ}{{\u{G}}}1
           {±}{{$\pm$}}1 {−}{{$-$}}1 {Δ}{{$\Delta$}}1}
\lstdefinestyle{kabuk}{language=bash,backgroundcolor=\color{black!85},basicstyle=\ttfamily\footnotesize\color{white},
  commentstyle=\color{green!70},keywordstyle=\color{white},stringstyle=\color{yellow!80}}

\newcommand{\OC}{O$-$C}
\newcommand{\Pdot}{\ensuremath{\dot{P}}}
\newcommand{\msyil}{\,ms\,yıl$^{-1}$}

\begin{document}

% ===================================================== KAPAK
\begin{titlepage}
\centering
\vspace*{1.5cm}
{\LARGE\bfseries Sıcak Jüpiterlerde Yörünge Dönemi Değişimi\par}
\vspace{0.4cm}
{\Large O$-$C Diyagramlarıyla Adım Adım Proje Kılavuzu\par}
\vspace{0.3cm}
{\large (Sıfırdan, “mala anlatır gibi”)\par}
\vspace{1.2cm}
{\large İdris Avlar\par}
{Atatürk Üniversitesi, Astronomi ve Uzay Bilimleri (lisans)\par}
\vspace{0.6cm}
{TÜBİTAK 2209-A Üniversite Öğrencileri Araştırma Projeleri — 2026 Çağrısı\par}
\vspace{1.5cm}
\begin{tcolorbox}[colback=white,colframe=black!60,width=0.9\textwidth]
\small
\textbf{Bu belge ne?} Hocanın verdiği görevin (14 aday $\rightarrow$ 3–4 hedef $\rightarrow$ O$-$C diyagramı
$\rightarrow$ dönem değişti mi? $\rightarrow$ değiştiyse hangi fiziksel etki?) baştan sona, hiçbir adımı atlamadan
anlatıldığı çalışma kılavuzudur. İçindeki bütün kodlar gerçek veriyle çalıştırılmış ve sonuçları
yazılmıştır. \textbf{Kısım I} TÜBİTAK formuna yapıştırılacak resmi metinlerdir; \textbf{Kısım II} sana özel kılavuzdur.\\[0.4em]
\textbf{Bu belge ne değil?} Eski “Detaylı Proje Raporu”nun yerine geçen, kapsamı \emph{daraltılmış} plandır.
Gaia astrometrisi, 26 adayın tamamı, teleskop gözlemi ve yeni tayf (RV) gözlemi \textbf{bu projede yoktur}.
\end{tcolorbox}
\vfill
{\small Kasım 2026}
\end{titlepage}

\tableofcontents
\newpage

% ===================================================== BÖLÜM 0
\section{Önce Bunu Oku: Proje Ne, Hoca Ne İstedi, Sen Ne Yapacaksın?}

\subsection{Hocanın söyledikleri (kelimesi kelimesine)}

\begin{hoca}[Hocanın direktifi]
“O makaledeki 26 kaynaktan LOOCV testini geçen 14 tanenin birkaç tanesini seç ama bir gerekçe sun bunun için.
Ayrıca bu testin ne olduğunu öğren; yani o testi geçip geçmemek ne anlama geliyor onu da öğren.
Bu seçtiğin gezegenlerin O$-$C diyagramlarını çiz. Yörünge dönemlerinde değişim var mı yok mu kontrol edeceğiz.
TESS ve VarAstro gibi database/katalog verilerini dahil edeceksin, literatürdeki verileri yani.
Yörünge dönemi değişimi varsa: 1.~Tidal interaction (effect), 2.~Apsidal motion, 3.~Rømer effect,
4.~Applegate effect, 5.~3.~cisim, 6.~Mass loss — bunlardan hangisinden etkilendiğini tespit etmen gerekiyor.”
\end{hoca}

Kağıttaki notlar da aynı şeyi söylüyor: “26 kaynak $\rightarrow$ pass $\rightarrow$ kullan; 14 kaynak $\rightarrow$
O$-$C'lerini çiz; yörünge dönemlerinde değişim var mı yok mu kontrol edeceğiz; TESS, katalog verileri dahil edilecek
(database: VarAstro gibi, literatürdeki veriler); yörünge dönem değişimi var ise $\rightarrow$ Tidal, Apsidal,
Rømer, Applegate, 3.~cisim, Mass loss.” İkinci kağıtta O$-$C diyagramının eksenleri çizilmiş:
\textbf{dikey eksen O$-$C, yatay eksen zaman ya da epok}. Bu belgenin tamamı bu iki kağıdın açılımıdır.

\subsection{Bu görevin altı adımı (tek bakışta)}

\begin{enumerate}[label=\textbf{Adım \arabic*:},leftmargin=*]
  \item \textbf{LOOCV'yi öğren} (Bölüm~\ref{sec:loocv}). Hocanın sana \emph{kişisel olarak} öğrenmeni istediği konu;
        projenin yöntemi değil. Testin ne yaptığını ve geçip geçmemenin anlamını kendi cümlelerinle anlatabilecek hale gel.
        Projede yalnızca seçimin başlangıç listesi olarak geçer (“Wang vd.'nin LOOCV'yi geçen 14 adayı”).
  \item \textbf{14 adaydan birkaçını gerekçeyle seç} (Bölüm~\ref{sec:secim}). Seçim bu belgede yapıldı:
        önceden yazılmış üç ölçütten kendiliğinden \textbf{WASP-12\,b, WASP-4\,b, HAT-P-7\,b} ve \textbf{TrES-1\,b}
        çıkıyor. Gerekçe, sayılarla birlikte hazır.
  \item \textbf{Geçiş zamanlarını topla} (Bölüm~\ref{sec:oc}). Hazır kataloglardan (Ivshina \& Winn 2022,
        Wang vd. 2024, ETD/VarAstro, ExoClock) ve TESS'ten. \emph{Gözlem yapmayacaksın.}
  \item \textbf{O$-$C diyagramını çiz ve dönem değişimini ölç} (Bölüm~\ref{sec:oc} ve \ref{sec:kod}).
        Doğrusal ve kuadratik model, $\dot P$, $\Delta$BIC, 3$\sigma$.
  \item \textbf{Değişim varsa sebebini bul} (Bölüm~\ref{sec:etkiler}, \ref{sec:rv}, \ref{sec:karar}).
        Altı etkinin her birini, sırayla, basit testlerle eleyeceksin. Gelgit ile Rømer'i ayırmak için
        \emph{arşivdeki} hazır radyal hız (RV) tablolarına bakacaksın (tayfçeker gerekmez).
  \item \textbf{Formu doldur ve gönder} (Kısım I ve Bölüm~\ref{sec:basvuru}).
\end{enumerate}

\subsection{Kapsamda olan / olmayan (hocanın sözünden çıkmamak için)}

\begin{center}
\begin{tabularx}{\textwidth}{>{\raggedright\arraybackslash}X >{\raggedright\arraybackslash}X}
\toprule
\textbf{KAPSAMDA (yapılacak)} & \textbf{KAPSAM DIŞI (yapılmayacak)}\\
\midrule
Wang vd. (2024)'nin LOOCV'yi geçen 14 adayından 4'ünün seçimi & 26 adayın tamamının yeniden analizi\\
Katalog/literatür geçiş zamanları (I\&W 2022, Wang 2024, ETD/VarAstro, ExoClock, makaleler) & Teleskopla yeni geçiş gözlemi (ATA50/T100 vb.)\\
TESS geçiş zamanları (hazır tablolar + Wang sonrası yeni TESS sektörleri) & Yeni tayf/RV gözlemi (tayfçekerimiz yok)\\
O$-$C diyagramı, doğrusal/kuadratik model, BIC & Gaia/Hipparcos astrometrisi, Gaia DR4; LOOCV'yi kendi verimize uygulamak\\
6 fiziksel etkinin basit testlerle değerlendirilmesi & Q$'_\star$ için Monte Carlo, yıldız evrim modelleri\\
Arşivdeki \emph{hazır} RV tablolarından eğilim okumak (yalnızca Rømer testi için) & Karmaşık LTTE/apsis MCMC modelleri (yalnızca gerekirse, basit sürüm)\\
\bottomrule
\end{tabularx}
\end{center}

\begin{dikkat}[RV konusu: tayfçeker yok, sorun da yok]
Bu projede \textbf{hiçbir tayf gözlemi yapılmayacak}. Radyal hız ölçümleri yıllar önce başka ekipler tarafından
(Keck/HIRES, HARPS, HARPS-N, SOPHIE, CORALIE ile) yapılmış ve \textbf{DACE arşivinde ve VizieR'de} herkese açık yayımlanmıştır.
Sen yalnızca bu hazır sayıları indirip “yıldızın hızı yıllar içinde düzenli olarak artıyor/azalıyor mu?” diye
bakacaksın. Bu, bir Excel tablosundaki sayılara doğru uydurmakla aynı zorlukta bir iştir. Hatta Bouma vd.
(2020)'nin Tablo~5'i bu hesabı senin hedeflerin için \emph{zaten yapmış}: doğrudan okuyup O$-$C sonucunla
karşılaştırman yeterli (Bölüm~\ref{sec:rv}).
\end{dikkat}

\subsection{Şu an ne yapacaksın? (Önce başvuru, sonra kod)}

\begin{basit}[SIRA (11 Kasım 2026, 17:30'a kadar)]
2209-A formu ön çalışmayı \emph{zorunlu} tutmaz (“yapıldıysa sunulması beklenir”). Ön çalışma şekilleri bu belgede
zaten hazır. Bu yüzden \textbf{önce başvuru}, kod sonra.
\begin{enumerate}[leftmargin=*]
\item \textbf{Bugün}: Hocana seçimi ve gerekçesini göster (Bölüm~\ref{sec:secim}'teki hazır paragraf ve Kısım I Tablo~1). Onay al.
\item \textbf{Form}: Kısım I'deki metinleri Word formuna yapıştır; Şekil~1, Şekil~2, Tablo~1, Tablo~2'yi Bölüm~5'e koy.
      Danışman adını yaz.
\item \textbf{Hoca kontrolü}: Formu hocana gönder, düzeltmeleri işle. Ekler hariç en fazla 20 sayfa, Arial 9.
\item \textbf{Gönder}: En geç \textbf{10 Kasım} akşamı TÜBİTAK sistemine yükle; danışmana onay için bir gün bırak.
\item \textbf{Sonra (başvurudan sonra ya da zaman kalırsa)}: Bölüm~\ref{sec:kod}'daki kodu Colab'da bir kez çalıştır.
      Bu başvuru için değil, hoca ya da jüri “bunu nasıl yaptın?” diye sorduğunda anlatabilmen için gerekli.
\end{enumerate}
Gün gün plan: Bölüm~\ref{sec:basvuru}.
\end{basit}

% ===================================================== BÖLÜM 1
\section{Sıfırdan Temel Kavramlar}
\label{sec:kavram}

Bu bölüm hiçbir şey bilmediğini varsayar. Her kavramın yanında “bu projede neden lazım?” sorusunun cevabı var.

\subsection{Sıcak Jüpiter ve geçiş (transit)}
\textbf{Sıcak Jüpiter}, kütlesi Jüpiter kadar olan ama yıldızına çok yakın (yörünge dönemi 1–5 gün) dolanan
gezegendir. Karşılaştırma için: Merkür 88 günde, bu gezegenler 1–3 günde bir tur atar.

\textbf{Geçiş (transit)}: Gezegen, bizim bakış doğrultumuzda yıldızın önünden geçerken yıldızın ışığının küçük bir
kısmını (\%0,5–2) keser. Yıldızın parlaklığını zamana karşı çizersen (\textbf{ışık eğrisi}) düzenli aralıklarla
tekrarlanan küçük “çukurlar” görürsün. Her çukur bir geçiştir.

\textbf{Orta geçiş zamanı ($T_c$, “mid-transit time”)}: Çukurun tam ortasının zamanı. Bu projede ölçtüğümüz,
topladığımız, çizdiğimiz \emph{tek} gözlemsel büyüklük budur. Örnek: WASP-12\,b'nin TESS'te ölçtüğüm bir geçişi
$T_c = 2460235{,}652614$ BJD$_\mathrm{TDB}$ $\pm$ 44 saniye.

\textbf{Yörünge dönemi ($P$)}: Art arda iki geçiş arasındaki süre. WASP-12\,b için $P \approx 1{,}0914$ gün
(26 saat 11 dakika).

\textbf{Epok ($E$ veya $N$)}: Bir başlangıç geçişinden bu yana geçen \emph{tam tur sayısı}. Başlangıç geçişi
$E=0$, bir sonraki $E=1$, bir önceki $E=-1$ \dots{} Epok hep tam sayıdır.

\subsection{Zaman nasıl yazılır: JD, BJD, TDB, BTJD}
\begin{itemize}[leftmargin=*]
\item \textbf{Jülyen Günü (JD)}: Gökbilimciler tarihleri gün/ay/yıl yerine tek bir sayıyla yazar: MÖ 4713'ten beri
geçen gün sayısı. Örnek: 1 Ocak 2024 öğlen $\approx$ JD 2460311,0. Kesirli kısım günün kesridir
(0,5 gün = 12 saat; 0,001 gün = 86,4 saniye).
\item \textbf{Barisentrik JD (BJD)}: Dünya Güneş etrafında dönerken yıldızdan gelen ışık bize bazen 8 dakika
erken, bazen 8 dakika geç ulaşır. Bu yüzden bütün zamanlar Güneş Sistemi'nin kütle merkezine (barisentere)
taşınır. Bu düzeltme yapılmış zamana BJD denir.
\item \textbf{TDB ve UTC}: Saatin “cinsi”. UTC artık saniyeler içerir; TDB düzgün akar.
2017'den beri \textbf{TDB $-$ UTC = 69,184 s}. Yani biri BJD$_\mathrm{UTC}$, diğeri BJD$_\mathrm{TDB}$ yazmışsa
aynı geçiş arasında \textbf{bir dakikadan fazla} fark çıkar. Dönem değişimi sinyali birkaç dakikalık olduğu için
bu fark \emph{sahte sinyal üretir}. Kural: \textbf{her zaman BJD$_\mathrm{TDB}$ kullan}.
\item \textbf{HJD}: Güneş merkezli (heliosentrik) eski bir sistem. BJD$_\mathrm{TDB}$'den $\sim$1 dakika farklı
olabilir; dönüştürmek gerekir (kodu Bölüm~\ref{sec:kod}'da).
\item \textbf{BTJD}: TESS'in kendi zaman sayısı: BTJD = BJD$_\mathrm{TDB}$ $-$ 2457000. Yani TESS dosyasındaki
\texttt{TIME} sütununa 2457000 eklersen BJD$_\mathrm{TDB}$ elde edersin.
\end{itemize}

\begin{dikkat}[Gerçek bir tuzak: HAT-P-7\,b]
Ivshina \& Winn (2022) veritabanında HAT-P-7\,b'nin 9 eski geçiş zamanı yalnızca “BJD” olarak etiketlenmiştir
(UTC mi TDB mi belli değil). Bu 9 zamanı dahil edince $\dot P = +22{,}6 \pm 5{,}1$\msyil, çıkarınca
$\dot P = +40{,}3 \pm 7{,}0$\msyil{} çıkıyor (kendi hesabım). \textbf{Tek bir zaman sistemi sorusu sonucu iki katına
çıkarabiliyor.} Bu yüzden her kaynağın zaman sistemini makalesinden kontrol etmek bu projenin en önemli “temizlik”
işidir.
\end{dikkat}

\subsection{Birimler: ms/yıl ne demek?}
Dönem değişim hızı $\dot P = dP/dt$ ile gösterilir. Boyutsuzdur (saniye/saniye) ama alışkanlıkla
\textbf{milisaniye/yıl} (ms\,yıl$^{-1}$) yazılır. $\dot P = -30$\msyil{} demek: “her yıl yörünge dönemi
30 milisaniye kısalıyor” demektir. Çok küçük görünür ama her turda birikir; WASP-12\,b'de 10 yılda geçişler
yaklaşık 8 dakika erken gelir (Bölüm~\ref{sec:oc}).

Dönüşüm: $\dot P\,[\mathrm{ms\,yıl^{-1}}] = \dot P\,[\text{boyutsuz}] \times 3{,}156\times 10^{10}$
(1 yıl = $3{,}156\times10^{7}$ s, 1 s = 1000 ms).

\subsection{Hata, sigma ($\sigma$) ve “3$\sigma$”}
Her ölçümün bir belirsizliği (hata çubuğu) vardır: $T_c = 2460235{,}652614 \pm 0{,}000508$ gün demek
“gerçek değer büyük olasılıkla bu aralıkta” demektir. Bir sonucun \textbf{3$\sigma$ anlamlı} olması, sonucun
sıfırdan en az 3 hata payı uzakta olması demektir: $|\dot P| / \sigma_{\dot P} \ge 3$. Örnek:
$\dot P = -28{,}6 \pm 1{,}0$ ise oran 28,6 $\Rightarrow$ çok anlamlı. $\dot P = -23 \pm 6{,}7$ ise oran 3,4
$\Rightarrow$ sınırda.

\subsection{En küçük kareler, $\chi^2$ ve BIC}
\textbf{Model uydurmak}: Noktalara bir eğri (doğru ya da parabol) geçirmek. En iyi eğri, noktalardan
\emph{hata çubuğuna göre ağırlıklı} uzaklıkların karelerinin toplamını en küçük yapan eğridir:
\[
\chi^2 = \sum_{i=1}^{n} \left( \frac{\text{gözlenen}_i - \text{model}_i}{\sigma_i} \right)^2 .
\]
$\chi^2$ küçükse model noktalara iyi uyuyor demektir.

\textbf{BIC (Bayes Bilgi Ölçütü)}: Daha çok parametreli model hep biraz daha iyi uyar; bu “haksız avantajı”
cezalandırmak için:
\[
\mathrm{BIC} = \chi^2 + k \ln n ,
\]
$k$: modeldeki serbest parametre sayısı (doğru için 2, parabol için 3), $n$: nokta sayısı.
\textbf{BIC'i küçük olan model kazanır.} $\Delta\mathrm{BIC} = \mathrm{BIC}_\text{doğrusal} -
\mathrm{BIC}_\text{kuadratik} > 10$ ise parabol “çok güçlü biçimde” tercih edilir (Kass \& Raftery 1995;
Wang vd. 2024 bu eşiği kullanır).

\subsection{Işık-yolu, radyal hız (RV) ve ivme}
\textbf{Radyal hız (RV)}: Yıldızın bize doğru ya da bizden uzağa hızı. Tayfçizgilerinin Doppler kaymasından
ölçülür. Gezegen yıldızı her turda biraz çeker, bu yüzden RV gezegenin dönemiyle bir sinüs çizer
(sıcak Jüpiterlerde genliği $\sim$200 m/s). Eğer sistemde uzakta \emph{ikinci bir ağır cisim} varsa, bütün sistem
yavaşça ivmelenir ve RV'de bu sinüsün üstüne yıllar boyunca süren \textbf{düz bir eğilim (trend)} eklenir:
$\dot\gamma = dv_r/dt$ (m\,s$^{-1}$\,yıl$^{-1}$). Bu eğilim, Rømer etkisinin “parmak izi”dir (Bölüm~\ref{sec:rv}).

